\chapter{Persistent Storage}

CharlotteOS treats storage as a layered stack of userspace services, each
speaking a well-defined protocol through endpoint IPC. No part of the storage
stack runs in the kernel — the kernel grants an NVMe driver an MMIO window
and an interrupt, and the rest is composed in userspace.

\section{The storage stack}

\begin{verbatim}
Filesystem ("fs")           ← charlotte-protocol-fs
    |
Object store ("obj")        ← charlotte-protocol-objstore
    |
Block driver ("blk0")       ← charlotte-protocol-block
    |
NVMe controller             ← MMIO + DMA + MSI-X (granted at boot)
\end{verbatim}

\section{NVMe block driver}

The NVMe driver is a userspace EL0 service implementing the NVM Express 1.4
base specification. It receives an \texttt{MmioRegion} capability covering the
controller's BAR0 register window and an \texttt{Interrupt} capability for I/O
completion signalling. It never names a physical address or interrupt vector
--- both arrive through the config-page contract.

\subsection{Initialisation}

The driver follows the standard NVMe initialisation sequence:

\begin{enumerate}
    \item \textbf{Controller reset:} disable CC.EN, wait for CSTS.RDY = 0.
    \item \textbf{Admin queue setup:} allocate physically contiguous memory for the
        Admin Submission Queue (ASQ, 32 entries) and Admin Completion Queue
        (ACQ, 32 entries). Write AQA, ASQ, and ACQ registers. Enable the
        controller and wait for CSTS.RDY = 1.
    \item \textbf{Identify controller:} submit an Identify command (CNS=1) through
        the admin SQ. Read capabilities, version, and supported features.
    \item \textbf{Set Features (Number of Queues):} declare how many I/O queue
        pairs the driver intends to use. QEMU's NVMe controller requires this
        before accepting I/O queue creation commands.
    \item \textbf{Identify namespace:} submit an Identify command (CNS=0, NSID=1)
        to read the namespace size (NSZE) and the active LBA format (FLBAS),
        which determines the block size.
    \item \textbf{Create I/O Completion Queue:} allocate memory, submit a Create
        I/O CQ admin command. The CDW11 register must set bit 0 (PC =
        Physically Contiguous) to 1 --- this is required by the NVMe
        specification and enforced by QEMU's device model.
    \item \textbf{Create I/O Submission Queue:} allocate memory, submit a Create
        I/O SQ admin command, associating it with the I/O CQ created in the
        previous step.
\end{enumerate}

After initialisation the driver registers \texttt{"blk0"} with the name service
and enters its main loop, blocking on \texttt{cq\_wait} for endpoint messages
and I/O completions.

\subsection{I/O path}

\textbf{OP\_READ} receives a BorrowWrite memory attachment from the caller.
The driver reads the physical address of the attached pages via
\texttt{memory\_get\_phys}, constructs an NVM Read command with the starting
LBA, block count, and PRP1 pointing to the caller's pages, and submits it to
the I/O SQ. The NVMe controller DMAs the data directly into the caller's
pages. When the completion arrives (polled via the I/O CQ), the driver replies
and the kernel revokes the borrow, returning the now-filled buffer to the
caller.

\textbf{OP\_WRITE} receives a Move memory attachment. The driver reads the
physical address, constructs an NVM Write command, and submits it. The
controller DMAs from the caller's pages. On completion the driver replies; the
buffer was moved and is consumed.

\textbf{OP\_FLUSH} submits an NVM Flush command to ensure all previously
written data reaches durable media before replying.

\subsection{Lessons from bringup}

The NVMe driver surfaced several non-obvious requirements:

\begin{itemize}
    \item \textbf{CDW11 bit 0 (PC) must be 1.} QEMU's NVMe device model
        requires physically contiguous queues. Without this bit, the Create
        I/O CQ command returns status 0x4002 regardless of all other
        parameters.
    \item \textbf{SQE layout is exacting.} The Command Identifier occupies
        bits 31:16 of DW0, the opcode occupies bits 7:0. The NSID shares
        DW0's high bits with DW0 in the first u64. PRP1 is at byte offset
        24, CDW10-11 at byte offset 40. Getting any of these wrong produces
        silent failures.
    \item \textbf{Doorbell stride must be derived from CAP.DSTRD.} On QEMU
        this is 0 (stride = 4 bytes), but the architecture supports strides
        up to $4 \ll 15$ bytes.
    \item \textbf{MMIO mapping must cover the doorbell region.} The NVMe
        controller registers occupy offsets 0x0000--0x003F within BAR0, but
        doorbell registers start at offset 0x1000. A single-page MMIO grant
        is insufficient --- at least two pages are needed.
    \item \textbf{On TCG, \texttt{cq\_wait} is required for VM exits.} A
        pure spin-loop prevents QEMU from processing I/O because the guest
        never yields. The driver must block on \texttt{cq\_wait} to let the
        host process NVMe commands.
\end{itemize}

\section{Block device protocol}

The block protocol (\texttt{charlotte-protocol-block}) is a minimal
\texttt{no\_std} crate defining the endpoint IPC interface between block
device consumers and drivers:

\begin{longtable}{|c|l|p{7cm}|}
\hline
\textbf{Opcode} & \textbf{Name} & \textbf{Description} \\
\hline
1 & OP\_INFO & Query block size and total block count. Scalar call/reply. \\
2 & OP\_READ & Read blocks into a BorrowWrite memory object. \\
3 & OP\_WRITE & Write blocks from a Move memory object. \\
4 & OP\_FLUSH & Flush write cache to durable media. \\
5 & OP\_TRIM & Discard a block range (optional). \\
\hline
\end{longtable}

The protocol is transport-agnostic: the same interface works for NVMe, AHCI,
virtio-blk, or a ramdisk. Consumers speak \texttt{OP\_READ}/\texttt{OP\_WRITE}
without knowing the underlying hardware.

\section{Object store}

The persistent object store (\texttt{charlotte-protocol-objstore}) provides
crash-safe, dynamically-sized blob storage on top of a block device. Objects
are identified by 64-bit monotonically-increasing IDs.

\subsection{On-disk layout}

\begin{verbatim}
LBA 0:     superblock (magic, version, generation, next_object_id)
LBA 1:     free block bitmap (1 bit per block)
LBA 2--33: object directory (512 entries × 32 bytes:
             id | flags | size_bytes | first_extent_lba)
LBA 34+:   data extents (linked list of blocks per object)
\end{verbatim}

\subsection{Crash safety}

New data extents are written to disk first, then the directory entry's
\texttt{first\_extent\_lba} pointer is updated atomically (a single-block
write). If a crash occurs between writing the extent and updating the
directory, the orphaned blocks are unreachable but the object reverts
consistently to its previous state. No journal, no double-writes --- the
pointer swap is the commit.

\subsection{Operations}

\begin{itemize}
    \item \textbf{CREATE} --- allocate a directory slot, return a new object ID.
    \item \textbf{WRITE} --- replace the object's content. The caller moves a
        memory object; the store writes it to a new extent and updates the
        directory pointer.
    \item \textbf{READ} --- return the object's data as a moved memory object.
    \item \textbf{DELETE} --- mark the directory entry as free, return the
        extent blocks to the free bitmap.
    \item \textbf{FLUSH} --- flush the underlying block device.
\end{itemize}

\section{Native filesystem}

The filesystem (\texttt{charlotte-protocol-fs}) builds hierarchical naming on
top of the object store. Directories and files are stored as objects; the root
directory is a fixed object ID (100).

\subsection{Directory encoding}

Each directory object stores entries as a packed sequence:

\begin{verbatim}
[name_len: u32 LE][name: UTF-8][file_id: u64 LE][flags: u32 LE][size: u64 LE]
\end{verbatim}

A \texttt{name\_len = 0} marker terminates the list. Flags distinguish
directories (bit 0 set) from files (bit 0 clear).

\subsection{Path resolution}

Paths are walked client-side by chaining \texttt{OP\_LOOKUP} calls. The
filesystem service does not understand path strings --- it only resolves
single-component lookups within a given directory. This keeps the filesystem
protocol minimal: six operations (LOOKUP, CREATE, READ, WRITE, DELETE, LIST)
plus FLUSH.

\subsection{Host inspection}

A Python tool (\texttt{scripts/fs-inspect.py}) reads the raw disk image and
walks the object store and filesystem on-disk structures. It provides tree
view, hex dump, cat, and format commands, enabling post-mortem debugging
without booting CharlotteOS:

\begin{verbatim}
python3 scripts/fs-inspect.py os-images/nvme-disk.img tree
python3 scripts/fs-inspect.py os-images/nvme-disk.img cat /raft/state
\end{verbatim}

\section{Why this matters for critical software}

\begin{enumerate}
    \item \textbf{Zero kernel code in the storage path.} A bug in the NVMe
        driver, the object store, or the filesystem corrupts that service's
        address space, not the kernel. The kernel grants MMIO and interrupts;
        everything else is userspace.

    \item \textbf{Transport independence.} The block protocol abstracts the
        hardware. The same object store and filesystem work over NVMe, AHCI,
        virtio-blk, or a network block device --- the kernel does not need to
        know which.

    \item \textbf{Crash-safe by construction.} The object store's atomic
        pointer-swap commit eliminates the need for journaling at the storage
        layer. A crash during a write leaves the previous version intact.

    \item \textbf{Capability-based access.} A client that holds a connection to
        \texttt{"blk0"} can read and write blocks. A client that does not hold
        that connection cannot. There is no ambient filesystem namespace to
        escape.

    \item \textbf{Host-inspectable.} The disk image is a raw file that can be
        examined with standard tools and the included Python inspector. No
        proprietary on-disk format, no opaque blobs.
\end{enumerate}

\endinput
